\documentclass[10pt,a4paper]{amsart}
\usepackage[margin=19mm]{geometry}
\usepackage{amsmath,amssymb,mathtools}
\usepackage[scr=rsfs,cal=euler]{mathalfa}
\usepackage{xfrac}
\usepackage[unicode,hidelinks,bookmarks=false]{hyperref}
\newcommand{\ie}{i.e.\ }
\newcommand{\cf}{cf.\ }
\newcommand{\resp}{respectively\ }
\newcommand{\Subsection}{\subsection}
\providecommand{\nobreakdash}{\nobreak\hbox{-}}
\setlength{\emergencystretch}{2em}
\multlinegap0pt
\usepackage{bm}
\usepackage{bbm}
\usepackage{dsfont}
\usepackage{xspace}
\usepackage{cancel}
\usepackage{mathtools}
\usepackage[shortlabels]{enumitem}
\usepackage{amsthm}
\makeatletter
\@ifundefined{theorem}{\newtheorem{theorem}{Theorem}}{}
\@ifundefined{lemma}{\newtheorem{lemma}[theorem]{Lemma}}{}
\@ifundefined{proposition}{\newtheorem{proposition}[theorem]{Proposition}}{}
\makeatother
\newcommand{\R}{\mathbb{R}}
\newcommand{\g}{\mathfrak}
\renewcommand{\H}{\mathbb{H}}
\title[Polar homogeneous foliations on symmetric spaces of rank one]{Polar homogeneous foliations on symmetric spaces of rank one}
\author{Jos\'e Carlos D\'iaz-Ramos, Juan Manuel Lorenzo-Naveiro}
\date{Source: arXiv:2606.27912v1 (CC BY 4.0)}
\begin{document}
\maketitle

\noindent\emph{Source and licence.} Polar homogeneous foliations on symmetric spaces of rank one. Version arXiv:2606.27912v1 (CC BY 4.0), distributed under the Creative Commons Attribution 4.0 International licence, as recorded in the arXiv licence metadata for this exact version and shown on the abstract page https://arxiv.org/abs/2606.27912v1. Full rights notice: licenses/arxiv-2606.27912-CC-BY-4.0.txt.\par\smallskip

\noindent\textit{Editorial scope.} Counted unit: the Lemma whose source label is \texttt{lemma:spperp} in the article (source0.tex lines 1321--1336, source label \texttt{lemma:spperp}), delivered together with the article's own complete original proof of it (source0.tex lines 1338--1455, one \texttt{proof} environment of 118 lines). No mathematics is written, repaired or re-derived: statement and proof are the article's own text line for line. Required context retained uncounted: th:criterion (lines 739--752); lemma:B+galpha (lines 1041--1043). Every definition, display and label of those lines is retained unchanged. Omitted from this package and not redistributed: 1--738, 753--1040, 1044--1320, 1337--1337, 1456--1946. Nothing inside a retained excerpt is omitted or altered: every retained body is delivered exactly as the article's lines are written, so no omission receipt of any kind accompanies this package. Citations: the retained excerpts carry no citation command at all, so no bibliography entry of the article is transcribed and no reference is redistributed. Reference adaptation: none was needed and none was made; every reference inside the delivered unit resolves inside it, so no unresolved reference or citation is produced. Printed slips kept literal: no printed word, symbol, hypothesis or number of the counted statement or of its proof was altered. Layout and class: the article's own class is \texttt{amsart} with packages including geometry, amsthm, amssymb, enumerate, xcolor, graphics, hyperref, array, booktabs; the wrapper is the lane's pinned \texttt{amsart} editorial wrapper at 10pt on A4, which restates the theorem-like environments the retained text needs and adds the article's own macro definitions that the retained text needs (\texttt{\textbackslash H}, \texttt{\textbackslash R}, \texttt{\textbackslash g}), with extra wrapper packages (bm, bbm, dsfont, xspace, cancel, mathtools, enumitem). The delivered numbering is the unit's own. Assets: no retained span contains a figure environment, a graphics-inclusion command, a PDF-page inclusion, a table or a TikZ picture; no image file is copied into this package, the package declares no asset dependency and no image file is redistributed with it. Licence: version arXiv:2606.27912v1 of this article is distributed under the Creative Commons Attribution 4.0 International licence, as recorded in the arXiv licence metadata for this exact version and shown on the abstract page https://arxiv.org/abs/2606.27912v1. A full rights notice is delivered at licenses/arxiv-2606.27912-CC-BY-4.0.txt. Characters: every retained excerpt is pure ASCII, so the delivered engine, XeLaTeX with its default Latin Modern text font, reports no missing character.
\begin{proposition}\label{th:criterion}
	Let $M=G/K$ be a Riemannian symmetric space of noncompact type with Cartan decomposition $\g{g}=\g{k}\oplus\g{p}$.
	Let $H\subseteq G$ be a closed subgroup acting on $M$ in such a way that its orbits form a foliation on $M$.
	We define
	\[
	\g{h}_\g{p}^\perp=\big\{\xi\in\g{p}:\langle\xi,X\rangle=0,\text{ for all $X\in\g{h}$}\big\}.
	\]
	Then:
	\begin{enumerate}[\rm (i)]
		\item The group $H$ acts polarly on $M$ if and only if $\g{h}_\g{p}^\perp$ is a Lie triple system in $\g{p}$ and $[\g{h}_\g{p}^\perp,\g{h}_\g{p}^\perp]$ is orthogonal to $\g{h}$.
		\item The group $H$ acts hyperpolarly on $M$ if and only if $\g{h}_\g{p}^\perp$ is an abelian subspace of $\g{p}$.
	\end{enumerate}
	Moreover, if $H_\g{p}^\perp$ denotes the subgroup of $G$ whose Lie algebra is $[\g{h}_\g{p}^\perp,\g{h}_\g{p}^\perp]\oplus\g{h}_{\g{p}}^{\perp}$, then $H_\g{p}^\perp\cdot o$ is a section of the action of $H$ on $M$.
\end{proposition}

\begin{lemma}\label{lemma:B+galpha}
    The action of $H$ is orbit equivalent to the action of a group $\hat{H}$ whose Lie algebra is contained in $\g{t}\oplus\g{a}\oplus\g{n}$ and satisfying the following condition: there is $\xi_\alpha\in\g{g}_\alpha$ such that $B+(1-\theta)\xi_\alpha\in\hat{\g{h}}_\g{p}^\perp$.
\end{lemma}

\begin{lemma}\label{lemma:spperp}
Assume $\dim(\g{h}_\g{p}^\perp)_{\g{p}_{2\alpha}}=1$.
We can write
\[
\g{h}_{\g{p}}^\perp
=\R(B+(1-\theta)\xi_\alpha)\oplus(1-\theta)\g{w}\oplus\R(1-\theta)(\eta_\alpha+\eta_{2\alpha}),
\]
where $\g{w}$ is a totally real subspace of $\g{g}_\alpha$, $\xi_\alpha$ and $\eta_\alpha$ are linearly independent vectors of $\g{g}_\alpha$, $\eta_{2\alpha}\in\g{g}_{2\alpha}$ with $\lvert\eta_{2\alpha}\rvert^2=2$, $\H\xi_\alpha+\H\eta_\alpha$ is orthogonal to $\H\g{w}$, 
$2[\xi_\alpha,\eta_\alpha]=-\eta_{2\alpha}$,
and
\[
\eta_\alpha
=\frac{\langle\xi_\alpha,\eta_\alpha\rangle}{\lvert\xi_\alpha\rvert^2}\xi_\alpha
-\frac{1}{\lvert\xi_\alpha\rvert^2}J_{\eta_{2\alpha}}\xi_\alpha.
\]
\end{lemma}

\begin{proof}
By Lemma~\ref{lemma:B+galpha} and the assumption $\dim(\g{h}_{\g{p}}^\perp)_{\g{p}_{2\alpha}}= 1$, it is clear that we can write
$\g{h}_{\g{p}}^\perp
=\R(B+(1-\theta)\xi_\alpha)\oplus(1-\theta)\g{w}\oplus\R(1-\theta)(\eta_\alpha+\eta_{2\alpha})$,
where $\g{w}\subseteq\g{g}_\alpha$, $\xi_\alpha$, $\eta_\alpha\in\g{g}_\alpha$ are orthogonal to $\g{w}$, $\eta_{2\alpha}\in\g{g}_{2\alpha}$, and $\lvert\eta_{2\alpha}\rvert^2=2$ (thus, $\lvert\eta_{2\alpha}\rvert_{AN}^2=1$).

Since $\g{h}_{\g{p}}^\perp$ is totally real, it is clear that so is $\g{w}$.
Moreover, for each $U\in\g{w}$ and $X\in\g{g}_{2\alpha}$ we have
\[
0
=\langle J_X(B+(1-\theta)\xi_\alpha),(1-\theta)U\rangle
=\Bigl\langle (1-\theta)\Bigl(\frac{1}{2}X+J_X\xi_\alpha\Bigr),(1-\theta)U\Bigr\rangle
=2\langle J_X\xi_\alpha,U\rangle,
\]
which implies $\H\xi_\alpha$ is orthogonal to $\g{w}$.
Similarly, since $J_X\eta_{2\alpha}\in\g{a}\oplus\g{g}_{2\alpha}$ we get
\[
0
=\langle J_X(1-\theta)(\eta_\alpha+\eta_{2\alpha}),(1-\theta)U\rangle
=2\langle J_X\eta_\alpha,U\rangle,
\]
and we also obtain that $\H\eta_\alpha$ is orthogonal to $\g{w}$.
Finally,
\begin{align*}
0
&{}=\langle J_{X}(B+(1-\theta)\xi_\alpha),(1-\theta)(\eta_\alpha+\eta_{2\alpha})\rangle\\
&{}=\Bigl\langle (1-\theta)\Bigl(\frac{1}{2}X+J_{X}\xi_\alpha\Bigr),
(1-\theta)(\eta_\alpha+\eta_{2\alpha})\Bigr\rangle\\
&{}=2\langle J_{X}\xi_\alpha,\eta_\alpha\rangle+\langle X,\eta_{2\alpha}\rangle
=\langle 2[\xi_\alpha,\eta_\alpha]+\eta_{2\alpha},X\rangle,
\end{align*}
so $2[\xi_\alpha,\eta_\alpha]=-\eta_{2\alpha}$.
In particular, $\xi_\alpha$ and $\eta_\alpha$ are linearly independent vectors.

Let $U\in\g{g}_\alpha\ominus\g{w}$.
The vector $-\langle U,\xi_\alpha\rangle B+U-\frac{1}{2}\langle U,\eta_\alpha\rangle\eta_{2\alpha}$
is orthogonal to $\g{h}_{\g{p}}^\perp$.
Thus, for each $U\in\g{g}_\alpha\ominus\g{w}$ there exists $T_U\in\g{t}$ such that
\begin{equation}\label{eq:tangent-vectors}
T_U-\langle U,\xi_\alpha\rangle B+U-\frac{1}{2}\langle U,\eta_\alpha\rangle\eta_{2\alpha}\in\g{h}.
\end{equation}
Taking $U=\xi_\alpha$ and $U=\eta_\alpha$ in the previous expression, and then bracketing the result gives
\begin{align*}
&
\Bigl[T_{\xi_\alpha}-\lvert\xi_\alpha\rvert^2 B+\xi_\alpha
-\frac{1}{2}\langle \xi_\alpha,\eta_\alpha\rangle\eta_{2\alpha},\,
T_{\eta_\alpha}-\langle \eta_\alpha,\xi_\alpha\rangle B+\eta_\alpha
-\frac{1}{2}\lvert\eta_\alpha\rvert^2\eta_{2\alpha}\Bigr]\\
&\qquad{}=\frac{1}{2}\langle\xi_\alpha,\eta_\alpha\rangle\xi_\alpha
-\frac{1}{2}\lvert\xi_\alpha\rvert^2\eta_\alpha
+[T_{\xi_\alpha},\eta_\alpha]-[T_{\eta_\alpha},\xi_\alpha]\\
&\qquad\phantom{{}=}{}+\frac{1}{2}\bigl(\lvert\xi_\alpha\wedge\eta_\alpha\rvert^2-1\bigr)\eta_{2\alpha}
-\frac{1}{2}\lvert\eta_\alpha\rvert^2[T_{\xi_\alpha},\eta_{2\alpha}]
+\frac{1}{2}\langle\xi_\alpha,\eta_\alpha\rangle[T_{\eta_\alpha},\eta_{2\alpha}],
\end{align*}
which is an element of $\g{h}$.
Hence, taking inner product with the vectors $B+(1-\theta)\xi_\alpha$ and $(1-\theta)(\eta_\alpha+\eta_{2\alpha})\in\g{h}_{\g{p}}^\perp$ we deduce that
%\begin{align*}
%0 &{}=\langle[T_{\xi_\alpha},\eta_\alpha],\xi_\alpha\rangle,\\
%0 &{}=\frac{1}{2}\lvert\xi_\alpha\wedge\eta_\alpha\rvert^2-1
%-\langle[T_{\eta_\alpha},\xi_\alpha],\eta_\alpha\rangle,
%\end{align*}
%respectively.
%This implies
\begin{align}\label{eq:Txieta}
\langle[T_{\xi_\alpha},\eta_\alpha],\xi_\alpha\rangle
&{}=0,&
\langle[T_{\eta_\alpha},\xi_\alpha],\eta_\alpha\rangle
&{}=\frac{1}{2}\lvert\xi_\alpha\wedge\eta_\alpha\rvert^2-1.
\end{align}

We also calculate
\begin{equation}\label{eq:bracket-spperp}
\begin{aligned}
&{}[B+(1-\theta)\xi_\alpha,(1-\theta)(\eta_\alpha+\eta_{2\alpha})]\\
&\qquad{}=(1+\theta)\Bigl(\frac{1}{2}\eta_\alpha+\eta_{2\alpha}
+[\xi_\alpha,\eta_\alpha]-[\theta\xi_\alpha,\eta_\alpha]
-[\theta\xi_\alpha,\eta_{2\alpha}]\Bigr)\\
&\qquad{}=(1+\theta)\Bigl(-[\theta\xi_\alpha,\eta_\alpha]
+\frac{1}{2}\eta_\alpha+J_{\eta_{2\alpha}}\xi_\alpha+\frac{1}{2}\eta_{2\alpha}\Bigr).
\end{aligned}
\end{equation}
Since $\g{h}$ is orthogonal to $[\g{h}_{\g{p}}^\perp,\g{h}_{\g{p}}^\perp]$ by Proposition~\ref{th:criterion}, taking inner product with the element $T_{\eta_\alpha}-\langle\xi_\alpha,\eta_\alpha\rangle B+\eta_\alpha-\frac{1}{2}\lvert\eta_\alpha\rvert^2\eta_{2\alpha}\in\g{h}$,
and using $\langle J_{\eta_{2\alpha}}\xi_\alpha,\eta_\alpha\rangle=\langle[\xi_\alpha,\eta_\alpha],\eta_{2\alpha}\rangle=-1$ and~\eqref{eq:Txieta}, gives
\begin{align*}
0&{}=-\langle T_{\eta_\alpha},(1+\theta)[\theta\xi_\alpha,\eta_\alpha]\rangle
+\frac{1}{2}\lvert\eta_\alpha\rvert^2
+\langle J_{\eta_{2\alpha}}\xi_\alpha,\eta_\alpha\rangle
-\frac{1}{2}\lvert\eta_{\alpha}\rvert^2\\
&{}=-2\langle[T_{\eta_{\alpha}},\xi_\alpha],\eta_\alpha\rangle-1
=-\lvert\xi_\alpha\wedge\eta_\alpha\rvert^2+1.
\end{align*}
Using these results and 
$\lvert J_{\eta_{2\alpha}}\xi_\alpha\rvert^2 =\frac{1}{2}\lvert\eta_{2\alpha}\rvert^2\lvert\xi_\alpha\rvert^2 =\lvert\xi_\alpha\rvert^2$ we calculate
\begin{align*}
&\lvert\eta_\alpha\wedge(\langle\xi_\alpha,\eta_\alpha\rangle\xi_\alpha-J_{\eta_{2\alpha}}\xi_\alpha)\lvert^2\\
&\quad{}=\lvert\eta_\alpha\rvert^2
\lvert\langle\xi_\alpha,\eta_\alpha\rangle\xi_\alpha-J_{\eta_{2\alpha}}\xi_\alpha\rvert^2
-\langle\eta_\alpha,\langle\xi_\alpha,\eta_\alpha\rangle\xi_\alpha-J_{\eta_{2\alpha}}\xi_\alpha\rangle^2\\
&\quad{}=\lvert\eta_\alpha\rvert^2
(\langle\xi_\alpha,\eta_\alpha\rangle^2\lvert\xi_\alpha\rvert^2+\lvert J_{\eta_{2\alpha}}\xi_\alpha\rvert^2)
-(\langle\xi_\alpha,\eta_\alpha\rangle^2-\langle J_{\eta_{2\alpha}}\xi_\alpha,\eta_\alpha\rangle)^2\\
&\quad{}=\lvert\eta_\alpha\rvert^2
(\langle\xi_\alpha,\eta_\alpha\rangle^2\lvert\xi_\alpha\rvert^2+\lvert \xi_\alpha\rvert^2)
-(\langle\xi_\alpha,\eta_\alpha\rangle^2+1)^2
=(\lvert\xi_\alpha\wedge\eta_\alpha\rvert^2-1)(\langle\xi_\alpha,\eta_\alpha\rangle^2+1)=0.
\end{align*}
This implies that $\eta_\alpha$ and $\langle\xi_\alpha,\eta_\alpha\rangle\xi_\alpha-J_{\eta_{2\alpha}}\xi_\alpha$ are collinear, so we may write
\begin{equation*}
\eta_\alpha=\mu(\langle\xi_\alpha,\eta_\alpha\rangle\xi_\alpha-J_{\eta_{2\alpha}}\xi_\alpha), \quad \mu\in \R.
\end{equation*}
Taking inner product with $\eta_\alpha$ and recalling that
$\lvert\xi_\alpha\wedge\eta_\alpha\rvert=1=-\langle J_{\eta_{2\alpha}}\xi_\alpha,\eta_\alpha\rangle$,
the formula for $\eta_\alpha$ readily follows.\qedhere
%the formula for $\eta_\alpha$ follows readily from the fact that 
%$\lvert\langle\xi_\alpha,\eta_\alpha\rangle\xi_\alpha-J_{\eta_{2\alpha}}\xi_\alpha\rvert^2
%=\lvert\xi_\alpha\rvert^2\lvert\eta_\alpha\rvert^2$.
\end{proof}

\end{document}
